\section{Typed Shallow Embedding}
\label{contrib:tagless}

The functional conversion implemented in~\cite{engel2023translation} operates on an untyped version of \mk even though the target languages are \haskell and \ocaml. 
These languages feature expressive static type systems, and it would be a shame to forgo their benefits. 
Besides explicit passing of generators this approach also erases all information about the type of specific variables, and combine all constructors in one type called \lstinline{Term}. 
Figure \ref{fig:generatedterm} shows an example of such type, which integrates natural numbers, lists, and binary trees.
This makes values such as \lstinline{S (Cons (Leaf) (S Nil))} possible, even though they have no meaningful interpretation.  

\begin{figure}[h]
    \centering
    \includelines{code/BalanceoOI.hs}{2}{8} 
    \caption{Generated \emph{Term} type}
    \label{fig:generatedterm}
\end{figure}

To overcome these issues, we need a typed \mk embedding capable of linking each logical value to its ground counterpart. 
Existing candidates fall short.
\haskell embedding typedKanren~\cite{kudasov2024typedkanren} supplies the required logical‑type machinery but only for execution, not for inspecting or analysing relations. 
Simultaneously, \ocanren offers a typed \ocaml embedding yet targets a language less convenient for our conversion pipeline. 
We therefore develop a new, tagless‑final~\cite{carette2009finally} embedding in Haskell that lets programmers write type‑safe and composable relations, provides the correspondence between logic and ground types, and remains open to future extensions. 

The embedding consists of two interfaces expressed as \haskell typeclasses:  

\begin{itemize}
    \item \lstinline{LogicType} that builds correspondence between the logical and the underlying type.
    \item \lstinline{Kanren} that describes a way to interpret a \mk relation. 
\end{itemize}    

These typeclasses are then instantiated to provide concrete logic types and interpreters of relations. 

\subsection{LogicType}
\label{tagless:logic}

\mk interpreters do not operate on concrete \haskell types; instead they expect logical types whose structures mirror the original but contain placeholders for variables, also called \emph{holes}.
The \lstinline{LogicType} typeclass with associated type \lstinline{WithLogic} enforce the shape of a logical type and establish a correspondence between ground and logical types. 
To support multiple interpreters, logic types require variables to represent holes and have to be polymorphic over them, while a variable itself should be polymorphic over the type of its content.
As such, we get the following definition of \lstinline{LogicType} and \lstinline{WithLogic}. 
\includelines{code/Logic.hs}{12}{13}

We also introduce two helper types: \lstinline{Logic}, which represents a logical hole at top-level, and \lstinline{Var}---a variable that contains a logical type.

\includelines{code/Logic.hs}{9}{10}

To illustrate the interface, we will examine it step by step and implement it for a simple case of Peano numbers.
Given the usual implementation of a natural number \lstinline{Nat}, we can define the logic type in the following way. 
\includelines{code/LogicNat.hs}{0}{3}

To map the underlying type to its logical counterpart, \lstinline{LogicType} features functions \lstinline{project} and \lstinline{reify}.

\includelines{code/Logic.hs}{14}{15}

The function \lstinline{project} maps the value of an underlying type to a fully-ground logical value of a logical type, while \lstinline{reify} converts a logical value to its underlying value and returns \lstinline{Nothing} if the input is not fully-ground.
These functions should obey the following laws:

\includelines{code/Logic.hs}{28}{29}

A simple instance for the Peano natural numbers might look like this:
\includelines{code/LogicNat.hs}{15}{20}

To be able to inspect logical values, we provide the following representations of term-style construction. 
It mirrors the typical structure of \mk values: a constructor with a fixed number of arguments of a logical type. 

\includelines{code/Logic.hs}{1}{8}

\lstinline{Field} is a type-erasing GADT, which ensures arguments may have different types. 
\lstinline{Constructor} represents a particular constructor of a term, with an explicit name and a way to create a value.

\lstinline{LogicType} introduces the function \lstinline{quote} which obeys the following law. 
\includelines{code/Logic.hs}{24}{25}

Implementing \lstinline{quote} can be done with helper functions \lstinline{quote0}, \lstinline{quote1}, and others. 
For type-preserving translator, it is important for the names passed to the constructors in \lstinline{quote} to correspond to the constructor names of the underlying type, however it is not in general a requirement of the system.
\includelines{code/LogicNat.hs}{5}{6}

Finally, it should be possible to unify logical types; thus \lstinline{LogicType} contains the function \lstinline{unifyVal}. 
\includelines{code/Logic.hs}{16}{18}

In addition to values to be unified, this function also accepts a \emph{unification provider}.
Unification provider is supplied by the interpreter (see \ref{tagless:interpreters}) and conveys the meaning of the unification in  interpreter's domain.

% This function is in fact redundant in the typeclass --- it can be implemented entirely using \emph{quote}, and in fact such an implementation is provided. It is, however, included, as it is performance-critical to many typical miniKanren interpreters, and utilizing \emph{quote} for it would be a significant performance hit.

This function can be implemented for Peano numbers in the following way. 
\includelines{code/LogicNat.hs}{8}{10}

% Last two functions of the typeclass also could also be reimplemented using \emph{quote}, and included for performance: \emph{derefVal} and \emph{generate}.

The next function of \lstinline{LogicType} are \lstinline{derefVal} that is similar to \lstinline{reify}, but works on values that can contain variables.
To achieve this, it requires an environment that provides values for such variables through the first argument. 

\includelines{code/Logic.hs}{19}{21} 

Its implementation for natural numbers is as follows. 
\includelines{code/LogicNat.hs}{12}{13}

The goal of the last function, \lstinline{generate}, is to systematically enumerate all possible underlying values. 
It replaces the user-provided generators required in the untyped version \ref{contrib:translator}.

\includelines{code/Logic.hs}{22}{22}

For natural numbers, the generator is rather simple:

\includelines{code/LogicNat.hs}{22}{22}

Functions \lstinline{unifyVal}, \lstinline{derefVal}, and \lstinline{generate} are redundant in the typeclass and can be implemented only by using \lstinline{quote}.
Such an implementation is provided, however its performance may be insufficient for particular interpreters. 
Because of that, the functions are included in the interface and can have custom implementations. 


\subsection{Kanren}
\label{tagless:kanren}

The tagless-final approach represent a domain-specific language by an interface of polymorphic combinators. 
Their behavior is determined by interpreters---instances of the typeclass---rather than a fixed syntax tree resulting in extensible, type-safe embeddings. 

In our implementation such interface is called \lstinline{Kanren}.
A type that has a \lstinline{Kanren} instance denotes \emph{relations that return a value}.
Therefore, \lstinline{Kanren rel} means that \lstinline{rel Int} is a relation that produces a value of type \lstinline{Int}. 

Conjunctions are often expressed with a monadic bind, but bind enforces a fixed evaluation order---preventing the reordering of conjuncts that our functional conversion needs.
Instead, we require the \lstinline{Applicative} constraint which imposes no order. 
With this constraint it is impossible to use the value of an earlier conjunct in the following code; thus, for pure \mk relations, the type of the return value is fixed to unit \lstinline{()}. 
The purpose of this parameter is described in \ref{tagless:eval}.

Disjunction, by contrast, naturally fits the \lstinline{Alternative} interface: \lstinline{empty} represents failure, and \lstinline{<|>} models branching without ordering constraints. 
Thus, we simply require \lstinline{Alternative}, which matches the semantics of \mk choice exactly.

Finally, each \lstinline{Kanren} instance has to specify a \lstinline{KVar} type --- a variable associated with this relation. 
It does not need to contain a value, but it has to be functorial over the type it is designated to contain.
Thus, we have the following definition with basic \mk operations:

\includelines{code/Kanren.hs}{6}{9}

In our embedding we need two ways to introduce variables. 
The first one, which corresponds to $\Fresh_{v}$, introduces an empty variable. 
The other denotes an argument of a relation that will, upon introduction, be bound to arguments passed to the call. 
It is equivalent to inserting a unification after $\Fresh_{v}$, but is immediately visible whenever transformations of the relation are done.

\includelines{code/Kanren.hs}{1}{3}

The function \lstinline{fresh} therefore receives a description of the variable to create and a continuation that builds the rest of the relation using that variable, and returns the fully assembled relation: 

\includelines{code/Kanren.hs}{10}{10}

A handful of helpers and the \lstinline{ApplicativeDo} \haskell extension allows writing relations concisely: 
\includecode{code/AddoTF.hs}

\subsection{KanrenEval}
\label{tagless:eval}

Although the \lstinline{Kanren} typeclass makes it possible to define and interpret relations, it offers no mechanism to extract the result of relation execution. 
This may seem counterintuitive, but some interpreters, such as the normalizing interpreter, do not operate with the output, instead transforming the relation itself. 

Nevertheless, it is an important property for evaluators which is why we extend \lstinline{Kanren} with the \lstinline{KanrenEval} typeclass. 
\lstinline{KanrenEval} allows dereferencing variables which makes their value accessible as relation's return value.
An interpreter can then define how to extract the resulting value (see \ref{tagless:interpreters}).

\includelines{code/Kanren.hs}{12}{13}

Combined with \lstinline{derefVal} function from \lstinline{LogicType} (\ref{tagless:logic}), the complete evaluation function is straightforward: 

\includelines{code/Kanren.hs}{15}{17}

\subsection{Automatic Derivation of LogicType Instances}
\label{tagless:types}

\lstinline{LogicType} instances are highly constrained because of the requirement that the logic type mirrors \mk terms and remains interoperable with the underlying type. 
Only algebraic data types (and GADTs) qualify, therefore exactly one valid implementation of \lstinline{LogicType} is admitted per underlying type---a fact we exploit for automatic derivation. 

\lstinline{TemplateHaskell} is capable of generating this unique instance for a given algebraic data type. 
For performance reasons (see Section \ref{task:evaluation}) the implementations of functions in this typeclass must be annotated \lstinline{INLINABLE}. 
This way, the compiler statically resolves which function is to be called for any particular type.

\subsection{Examples of Interpreters}
\label{tagless:interpreters}

Any instance of \lstinline{Kanren} is called an interpreter, whether it evaluates a relation, transforms it, or does something completely different. 
In this subsection we will consider two examples of such interpreters. 

For instance, consider the standard \mk interpreter over a stream of substitutions. 
The result of its execution, or representation, is \lstinline{SubstKanren}. 
We add instances of \lstinline{Kanren} and \lstinline{KanrenEval} for the representation and provide the evaluation function \lstinline{runSubstKanren}. 
\includecode{code/SubstKanren.hs}

Note that the substitution is a partial function that relies on \lstinline{unsafeCoerce}, yet it is safe, as long as \lstinline{SVar} constructor is not used by the programmer. 
Two measures enforce this safety: a marker-parameter \lstinline{s} (inspired by the implementation of \lstinline{ST}), and the type tracking as part of a variable type. 
They ensure that the substitution always produces a value, and \lstinline{unsafeCoerce} is applied only between identical types.

Now let us consider a normalizing interpreter which does not evaluate a relation, instead transforming it. 
It might be encoded in the following way: 
\includecode{code/NormKanren.hs}

We omit the instances of \lstinline{Applicative} and \lstinline{Alternative} which contain the normalization logic, because they are not important to the overall structure. 

\subsection{Extensibility Example}

The major advantage of the tagless-final encoding is that additional structure can be introduced and composed directly in user code without any changes to the framework itself.
Consider normalization from the previous subsection that can transform any relation without changing its observable type.
\includelines{code/NormDemonstration.hs}{1}{9}


However, here normalization is totally transparent. 
Even though the structure of the relation is more constrained, it is not exposed and cannot be used in further transformations, because the only interface available is \lstinline{Kanren}.
We can combat this issue by defining a new typeclass, \lstinline{NormalizedKanren}, to expose the normalized shape and make it interoperable with the \lstinline{Kanren} typeclass.

\includecode{code/NormKanrenTypeclass.hs}

With this setup, any relation can be processed by a normalizing interpreter before being passed on to a conversion which will take full advantage of the normalized structure. 
Conversely, with an instance of \lstinline{NormalizedKanrenT}, any normalized relation can be restored to the \lstinline{Kanren} typeclass with the \lstinline{normalize} function, erasing the structure.
As a result, we demonstrated extensibility by providing two interoperable relation representations with no need to modify \lstinline{Kanren}, or any other core feature. 
